\documentclass[letterpaper, 10pt]{article}
\makeatletter
\providecommand*{\input@path}{}
\g@addto@macro\input@path{{../../}}
\makeatother
\usepackage[left=2cm, right=2cm, top=2cm, bottom=2cm]{geometry}
\usepackage{titling}
\setlength{\droptitle}{-5em}

% Packages not loaded by std_page.tex but required for article document class
\usepackage{amsmath}
\usepackage{amssymb}
\usepackage{mathtools}
\usepackage{graphicx}
\usepackage{xcolor}
\usepackage{longtable}
\usepackage{booktabs}
\usepackage{array}
\usepackage{hyperref}
\usepackage{framed}
\usepackage{colortbl}
\hypersetup{
    colorlinks=true,
    linkcolor=blue,
    filecolor=magenta,
    urlcolor=cyan,
}

% --- Color Definitions ---
\definecolor{reviewbg}{RGB}{245, 245, 255}       % light blue-grey for reviewer quotes
\definecolor{reviewtext}{RGB}{50, 50, 120}        % dark blue for reviewer quote text
\definecolor{revisioncolor}{RGB}{0, 100, 0}       % dark green for revision actions
\definecolor{rowsep}{RGB}{180, 180, 200}          % subtle row separator color
\definecolor{headerbg}{RGB}{230, 235, 245}        % header background

% --- Helper command for reviewer quotes ---
\newcommand{\rquote}[1]{\textcolor{reviewtext}{\emph{#1}}}

% --- Helper command for revision notes ---
\newcommand{\revision}[1]{\par\smallskip\textcolor{revisioncolor}{\textbf{Revision:}~#1}}

% Input workspace custom packages and commands
\input{../../ltx_core/ltx_pkgs.tex}
\newcommand{\bm}{}
\input{../../ltx_core/new_cmds.tex}

\newenvironment{reviewercomment}
{\begin{framed}\noindent\textbf{Reviewer Comment:}\itshape}
{\end{framed}}

\title{Author Responses to Peer Reviews \\ \large Joint MECC 2026 -- ASME JDSMC Submission \#26}
\author{Sesha Charla, Peter H. Meckl, Bin Yao}
\date{\today}

\begin{document}
\maketitle

We thank the Associate Editor and the reviewers for their constructive feedback and valuable suggestions on our
manuscript. We have carefully revised the manuscript to address all the points raised.

We are encouraged by the reviewers' recognition of the novelty and practical significance of our work. The Associate
Editor described it as addressing \emph{``an important and interesting research problem about modeling and parameter
    estimation of Diesel engine aftertreatment system with experimental validations.''} Reviewer~14 noted that \emph{``the
    proposed algorithm holds significant potential for practical engineering applications''} and that \emph{``the paper is
    well-structured, tightly integrating theoretical derivation with engineering validation.''} Reviewer~18 acknowledged
the practical value of the approach for \emph{``diagnosing catalyst aging in diesel Selective Catalytic Reduction (SCR)
    systems.''}

Below we provide responses to each of the specific comments. Concerns shared by multiple reviewers are
consolidated and addressed first.

\bigskip

\renewcommand{\arraystretch}{1.5}
\setlength{\tabcolsep}{6pt}
\begin{longtable}{
    >{\raggedright\arraybackslash\small}p{0.12\textwidth}
    >{\raggedright\arraybackslash\small}p{0.36\textwidth}
    >{\raggedright\arraybackslash\small}p{0.46\textwidth}}
    \toprule
    \rowcolor{headerbg}
    \textbf{Reviewers}         &
    \textbf{Reviewer Feedback} &
    \textbf{Rebuttal Comments and Revisions} \\ \midrule
    \endfirsthead
    \toprule
    \rowcolor{headerbg}
    \textbf{Reviewers}         &
    \textbf{Reviewer Feedback} &
    \textbf{Rebuttal Comments and Revisions} \\ \midrule
    \endhead
    \bottomrule
    \endfoot
    %==================================================================================================================
    % ROW 1: Necessary conditions / OBD feasibility (MAJOR — shared by AE + R14 + R18)
    %==================================================================================================================
    Associate Editor,\newline Reviewer~18,\newline Reviewer~14
                               &
    \rquote{AE: ``Justify the necessary condition of the proposed method based on sufficient data under vehicle working
        conditions.''}
    \par\medskip
    \rquote{R14: ``During real-world heavy-duty truck operations, the vehicle may encounter prolonged light-load
        cruising [\ldots] which might result in entire driving segments failing to trigger catalyst saturation at all. Under
        such scenarios, will the system generate false alarms?''}
    \par\medskip
    \rquote{R18: ``This is a severe bottleneck for real-world deployment, evidenced by the fact that the actual truck
        drive segments [\ldots] featured very few samples (often under 10\%) near saturation.''}
                               &
    Although only a fraction of the operating data is close to saturation (approximately 10\% in real-world truck data),
    the convex parameter estimation framework utilizes the \textbf{entire} dataset. The non-saturated data points
    establish the lower boundary constraints ($\eta(k) \le \eta_{\sat}(k)$), which restrict the parameter space and
    prevent the estimates from collapsing.
    \par\smallskip
    For on-vehicle implementation, catalyst utilization can be assessed in real time by monitoring readily available
    signals such as urea dosing rate, engine rpm, and the $\nox$ reduction estimate from sensors. When these indicate
    sustained low catalyst demand, the system defers its decision rather than executing the hypothesis test on
    insufficient data. \textbf{The detector can thus avoid producing false alarms under non-catalyst-saturating driving
        conditions}---the Wald test is simply not invoked.
    \revision{A dedicated paragraph discussing these practical OBD considerations, false-alarm prevention, and the impact
        on the in-use monitor performance ratio (IUPMR) has been added at the end of Section~6. A note on monitoring
        catalyst utilization via available signals has been added to the Limitations appendix (Appendix~B).}
    \\ \arrayrulecolor{rowsep}\midrule\arrayrulecolor{black}
    %==================================================================================================================
    % ROW 2: Threshold beta interpretation (MAJOR — R14)
    %==================================================================================================================
    Reviewer~14
                               &
    \rquote{``In the test-cell validation, the decision thresholds are set to $\beta_{\text{hFTP}} = 4.00$ and
        $\beta_{\text{RMC}} = 10.00$, respectively. Why is the threshold abruptly scaled up directly to
        $\beta_{\text{truck}} = 250$ in the real-world truck driving segment analysis?''}
                               &
    The threshold $\beta_{\text{truck}} = 250$ is significantly larger than the test-cell thresholds due to the
    finite-sample departure of $T_w$ from its asymptotic $\chi^2_p$ distribution under $\mathcal{H}_0$. Truck data has
    considerably fewer saturated-mode samples ($N_{\text{sat}} \approx 120$--$150$) compared to test-cell experiments
    ($N_{\text{sat}} \approx 200$--$500$), resulting in heavier tails in the null distribution. This is compounded by the
    less reliable estimation of the error variance under road conditions and the unaccounted variance of the reference
    parameter $\theta_{\text{sat}}^{dg}$.
    \revision{Two paragraphs have been added at the end of Section~6 justifying the threshold discrepancy and discussing
        the path toward statistically optimal threshold selection.}
    \\ \arrayrulecolor{rowsep}\midrule\arrayrulecolor{black}
    %==================================================================================================================
    % ROW 3: Independence of saturated mode / scope (MAJOR — R18)
    %==================================================================================================================
    Reviewer~18
                               &
    \rquote{``The saturated catalyst model's response does not depend on the previous state but only on the inputs,
        making the response independent of the initial conditions. This effectively means [\ldots] that the dynamics are
        irrelevant, and that is the reason a simple MLE point estimator suffices. In that sense, this paper does not fit the
        scope of MECC, because the system dynamics don't really matter.''}
                               &
    The goal of the referenced phrase was to convey that one mode (the saturated mode) of the switched dynamic model for
    $\nox$ reduction is independent of initial conditions. This property allows us to use the prediction error to
    classify data points as ``close to saturation'' when computing the statistics for aging detection. The next time-step
    $\nox$ reduction \textbf{still depends on current inputs} through a discrete-time difference equation whose output is
    governed by the sampling-to-residence-time ratio.
    \par\smallskip
    The saturated-mode model is not a static relationship---it is a discrete-time input-output model derived from molar
    conservation laws. The MLE estimator exploits the time-series structure of the saturated-mode dynamics. The
    non-saturated data points contribute dynamic inequality constraints. The overall framework is fundamentally a
    \textbf{dynamic systems identification and change detection problem}.
    \par\smallskip
    The requirement of catalyst saturation is a deliberate design choice that enables the core novelty: saturated
    operation simplifies the highly nonlinear, unobservable ammonia coverage dynamics into a static limit, allowing a
    convex parameter estimation framework \textbf{without requiring an on-board ammonia sensor}.
    \revision{Additional clarification has been added to the first paragraph of Section~4 to better convey the role of
        mode independence within the switched dynamic framework.}
    \\ \arrayrulecolor{rowsep}\midrule\arrayrulecolor{black}
    %==================================================================================================================
    % ROW 4: Significance of the paper (MODERATE — R12)
    %==================================================================================================================
    Reviewer~12
                               &
    \rquote{``The significance of this paper should be stated clearly.''}
                               &
    We have added a dedicated paragraph in the Introduction explicitly outlining the three main contributions:
    (1)~a first-principles discrete-time input-output model derived from molar conservation that captures the interplay
    between gas parcel residence time and sensor sampling intervals;
    (2)~a convex parameter estimation framework that bypasses the need for on-board $\ammonia$ sensors while remaining
    robust to the unidirectional bias of commercial $\nox$ sensors; and
    (3)~a Wald-test-based catalyst aging detector validated using both test-cell and real-world truck driving datasets.
    \revision{The significance paragraph has been added to the end of Section~1 (Introduction).}
    \\ \arrayrulecolor{rowsep}\midrule\arrayrulecolor{black}
    %==================================================================================================================
    % ROW 5: Experimental comparison — figures vs. tables (MODERATE — R12)
    %==================================================================================================================
    Reviewer~12
                               &
    \rquote{``The readers may want to see the clear comparison of experimental data and model results. It is better to
        present the figures but the numbers.''}
                               &
    The detection results span multiple trucks, multiple drive segments, and two
    data years. Presenting all results as individual time-series plots would obscure the cross-comparison that the tabular format enables. The existing tables (Tables~5--9) summarize the detection results for all trucks and drive segments in a compact format that facilitates direct comparison of aging detection outcomes.
    \revision{No additional figures were added. The tabular format was retained as the most effective presentation for
        the scope of the results.}
    \\ \arrayrulecolor{rowsep}\midrule\arrayrulecolor{black}
    %==================================================================================================================
    % ROW 6: Literature review (MINOR — R12)
    %==================================================================================================================
    Reviewer~12
                               &
    \rquote{``The literature review is not sufficient. The limitations of the existing works are not clear.''}
                               &
    The paper cites 24 references organized around three
    pillars: SCR modeling and control (9~references, including Hsieh \& Wang, Devarakonda et al., Ciardelli et al.),
    SCR aging diagnostics (6~references, including Ma \& Wang, Jain et al., Jiang et al.), and the underlying
    mathematical and statistical theory (10~references, including Kay, Wald, Trefethen). These references collectively
    justify the research objective, identify the gap, and map directly to the specific focus of model-based SCR aging
    diagnostics under catalyst saturation.
    \revision{Additional references to justify the $\beta$ discrepancy were added.}
    \\ \arrayrulecolor{rowsep}\midrule\arrayrulecolor{black}
    %==================================================================================================================
    % ROW 7: Formatting errors and writing issues (MINOR — AE + R12 + R14)
    %==================================================================================================================
    Associate Editor,\newline Reviewer~12,\newline Reviewer~14
                               &
    \rquote{AE: ``Some writing issues (e.g., Appendix-?? on page 5) need to be revised.''}
    \par\medskip
    \rquote{R12: ``There are some writing mistakes in this paper. Such as: the expression of `(Appendix-??)' is wrong.''}
    \par\medskip
    \rquote{R14: ``$\nox$, $\ammonia$, $\oxygen$, and $\nitrogen$ must be strictly rendered in roman font.''}
                               &
    1.~The writing issues were due to a missing cross-reference link caused by a LaTeX label mismatch while converting
    the paper from the JDSMC template to the MECC template. The broken cross-reference has been corrected and now
    compiles as (Appendix~B) in Section~3.
    \par\smallskip
    2.~All chemical symbols have been updated to roman font ($\nox$, $\ammonia$, $\oxygen$, $\nitrogen$) throughout the
    text, as suggested by Reviewer~14.
    \revision{Cross-reference fixed in Section~3; chemical symbols corrected throughout.}
    \\ \bottomrule
\end{longtable}

\end{document}
